% !TEX root = ../main.tex
\section{Problema 3: Secuencia Mágica}

\subsection{Descripción del Problema y Abstracción CSP}
% Explicar qué es una secuencia mágica de longitud n:
% Secuencia x_0, ..., x_{n-1} de enteros tales que x_i in 0..(n-1) y el número i aparece exactamente x_i veces.

\subsubsection{Parámetros y Datos de Entrada}
% Parámetro n (longitud de la secuencia).

\subsubsection{Variables de Decisión y Dominios}
% Variables x_0, ..., x_{n-1} con dominio 0..(n-1).

\subsubsection{Restricciones Básicas}
% Formalizar matemáticamente:
% Conteo: para cada i in 0..(n-1), sum_{j=0}^{n-1} (x_j == i) = x_i (o global_cardinality / count).

\subsection{Detalles de Implementación (\texttt{secuencia.mzn})}

\subsubsection{Restricciones Redundantes y Demostración Teórica}
% El taller exige evaluar e incluir dos restricciones redundantes:
% 1. \sum_{i=0}^{n-1} x_i = n
% 2. \sum_{i=0}^{n-1} (i - 1) x_i = 0  (es decir, (-1)*x_0 + ... + (n-2)*x_{n-1} = 0)
% Explicar matemáticamente por qué estas restricciones son redundantes y cómo ayudan a podar el espacio de búsqueda.

\subsubsection{Rompimiento de Simetrías}
% Analizar si existen simetrías en las soluciones del problema y si es posible imponer restricciones para romperlas.

\subsection{Estrategias de Búsqueda y Distribución}
% Anotaciones de búsqueda en MiniZinc para enumerar todas las secuencias mágicas.
% Comparar el comportamiento del solver con y sin las restricciones redundantes.

\subsection{Pruebas y Análisis Experimental}
% Pruebas variando la longitud n (ej. n = 4, 5, 10, 20, 50, etc.).
% Evaluar la reducción del árbol de búsqueda (nodos, fallos, tiempo).

\begin{table}[htbp]
    \centering
    \begin{tabular}{cccccc}
        \toprule
        $n$ & Modelo & Tiempo (s) & Nodos explorados & Nodos fallidos & Soluciones \\
        \midrule
        $10$ & Sin redundantes & -- & -- & -- & -- \\
        $10$ & Con redundantes & -- & -- & -- & -- \\
        $20$ & Sin redundantes & -- & -- & -- & -- \\
        $20$ & Con redundantes & -- & -- & -- & -- \\
        \bottomrule
    \end{tabular}
    \caption{Impacto de las restricciones redundantes en la Secuencia Mágica.}
    \label{tab:secuencia_resultados}
\end{table}

\subsection{Conclusiones del Ejercicio}
% Justificación de cómo las restricciones redundantes mejoran la propagación y reducen el espacio de búsqueda.
